% ---- EDy0/EDy1: Einteilung der elektromagnetischen Felder ----
\bereich[cKap0]{Einteilung der elektromagnetischen Felder}
\begin{formelreihe}
  \parbox[t]{\linewidth}{%
    \formelblock{Zeitlich konstant $\partial/\partial t=0$}{\begin{aligned}
      &\text{Elektrostatik: } \vec{H},\vec{B}=0\\[0.1em]%
      &\text{Magnetostatik: } \vec{E},\vec{D}=0\\[0.1em]%
      &\text{Stationaeres Stroemungsfeld}
    \end{aligned}}\\[0.25em]%
    \formelblock{Langsam zeitveraenderl.\ $\partial/\partial t\neq0$}{\begin{aligned}
      &\tfrac{\partial\vec{B}}{\partial t}\approx 0\text{: quasistat.\ el.\ Feld}\\[0.1em]%
      &\tfrac{\partial\vec{D}}{\partial t}\approx 0\text{: quasistat.\ mag.\ Feld}
    \end{aligned}}} &
  \parbox[t]{\linewidth}{%
    \formelblock{Schnell zeitveraenderlich}{\text{Elektromagnetisches Feld}}\\[0.25em]%
    \infoblock{%
      \textbf{Materialgleichungen:}\\
      $\vec{D}=\varepsilon\vec{E}$,\;
      $\varepsilon=\varepsilon_0\varepsilon_r$\\[0.15em]
      $\vec{B}=\mu\vec{H}$,\;
      $\mu=\mu_0\mu_r$\\[0.15em]
      $\varepsilon_0=8{,}854\cdot10^{-12}\,\mathrm{As/(Vm)}$\\
      $\mu_0=4\pi\cdot10^{-7}\,\mathrm{Vs/(Am)}$}} &
  \beispielblock{Lichtgeschwindigkeit}{c=\tfrac{1}{\sqrt{\varepsilon_0\mu_0}}
    \approx 3\cdot10^8\,\tfrac{\mathrm{m}}{\mathrm{s}}} \\
\end{formelreihe}
\bereichende
